\documentclass[11pt,a4paper]{article}

\usepackage[margin=1in]{geometry}
\usepackage{amsmath,amssymb,mathtools}
\usepackage{enumitem}
\usepackage{hyperref}

\hypersetup{
    colorlinks=true,
    linkcolor=black,
    urlcolor=blue
}

\newcommand{\Id}{\operatorname{Id}}
\newcommand{\Span}{\operatorname{span}}
\newcommand{\diag}{\operatorname{diag}}
\newcommand{\spec}{\sigma}

\title{Linear Algebra Problem Set 13}
\author{Taras Halay}
\date{}

\begin{document}

\maketitle

\section*{Exercise 5.15}

Since $A$ is orthogonal, assume $\lambda$ is an eigenvalue of $A$, then
\[
Ax=\lambda x
\]
for some $x$.

Then
\[
\|Ax\|=\|\lambda x\|=\|x\|
\]
by definition of orthogonal matrix.

Then
\[
\|\lambda x\|=|\lambda|\|x\|=\|x\|
\]
implies
\[
|\lambda|=1 \Rightarrow \lambda\in\{-1,1\}.
\]

So the only possible eigenvalues are $1$ and $-1$.

Since $A$ is diagonalizable, then there exists a matrix $P$ such that $P$ is invertible, and
\[
A=PDP^{-1}.
\]

As we know that $\det(A)=1$, hence
\[
\det(A)=\det(PDP^{-1})
\]
\[
=\det(P)\det(D)\det(P^{-1})
\]
\[
=\det(P)\det(D)\frac{1}{\det(P)}=\det(D).
\]

And since
\[
D=
\begin{pmatrix}
\lambda_1 & 0\\
0 & \lambda_2
\end{pmatrix},
\]
where $\lambda_1,\lambda_2$ are eigenvalues of $A$,
\[
\det(A)=\det(D)=\lambda_1\lambda_2=1.
\]

Since $\lambda_1,\lambda_2\in\{-1,1\}$ and $\lambda_1\lambda_2=1$, then
\[
D=I \quad \text{or} \quad D=-I.
\]

If $D=I$, then
\[
A=PIP^{-1}=PP^{-1}=I.
\]

If $D=-I$, then
\[
A=P(-I)P^{-1}=-PP^{-1}=-I.
\]

\section*{Exercise 5.6}

Since $L$ is an isomorphism, $L^{-1}$ exists.

Then consider
\[
M\circ L=L^{-1}\circ (L\circ M)\circ L.
\]

Then, let's compute the characteristic polynomial of $M\circ L$.

\[
\det(M\circ L-\lambda I)=
\det\bigl(L^{-1}\circ (L\circ M)\circ L-\lambda I\bigr)
\]
\[
=\det\bigl(L^{-1}\circ (L\circ M-\lambda I)\circ L\bigr)
\]
\[
=\det(L^{-1})\det(L\circ M-\lambda I)\det(L).
\]

Since $L$ is invertible,
\[
=\frac{1}{\det(L)}\det(L\circ M-\lambda I)\det(L)
=\det(L\circ M-\lambda I).
\]

Hence
\[
\det(L\circ M-\lambda I)=\det(M\circ L-\lambda I).
\]

\section*{Exercise 5.8}

Let
\[
A=\begin{pmatrix}
p & -q\\
q & p
\end{pmatrix},
\qquad p,q\in\mathbb R,\quad q\neq0.
\]

\[
\det(A-\lambda I)
=\det\begin{pmatrix}
p-\lambda & -q\\
q & p-\lambda
\end{pmatrix}
\]
\[
=(p-\lambda)^2+q^2>0
\]
over $\mathbb R$, since $q\neq0$.

So there are no real eigenvalues over $\mathbb R$.

Over $\mathbb C$,
\[
(p-\lambda)^2+q^2
=(p-\lambda-qi)(p-\lambda+qi)
\]
\[
=((p-qi)-\lambda)((p+qi)-\lambda).
\]

Hence
\[
\lambda_1=p-qi,\qquad \lambda_2=p+qi.
\]

Then $A$ has two distinct eigenvalues, so it is diagonalizable over $\mathbb C$.

For $\lambda_1=p-qi$,
\[
Ax=(p-qi)x
\]
means
\[
\begin{pmatrix}
p & -q\\
q & p
\end{pmatrix}
\begin{pmatrix}
x_1\\
x_2
\end{pmatrix}
=
(p-qi)
\begin{pmatrix}
x_1\\
x_2
\end{pmatrix}.
\]

So
\[
\begin{pmatrix}
px_1-qx_2\\
qx_1+px_2
\end{pmatrix}
=
\begin{pmatrix}
(p-qi)x_1\\
(p-qi)x_2
\end{pmatrix}.
\]

Thus
\[
qx_2=qix_1\Rightarrow x_2=ix_1,
\]
and
\[
qx_1=-qix_2\Rightarrow x_1=-ix_2.
\]

Therefore
\[
V^A_{\lambda_1}=\Span\left\{
\begin{pmatrix}
1\\
i
\end{pmatrix}
\right\}.
\]

For $\lambda_2=p+qi$,
\[
Ax=(p+qi)x
\]
gives us
\[
x_2=-ix_1,
\qquad
x_1=ix_2.
\]

Therefore
\[
V^A_{\lambda_2}=\Span\left\{
\begin{pmatrix}
1\\
-i
\end{pmatrix}
\right\}.
\]

\section*{Exercise 5.19}

\subsection*{(i)}

Let $\lambda\in\mathbb C$ be an eigenvalue of $L$. This means
\[
L(x_0,x_1,x_2,\ldots)=\lambda(x_0,x_1,x_2,\ldots),
\]
which means
\[
(x_1,x_2,\ldots)=(\lambda x_0,\lambda x_1,\ldots).
\]

Therefore
\[
x_{i+1}=\lambda x_i
\]
for all $i=0,\ldots,\infty$.

And this is exactly the definition of a complex infinite geometric sequence.

And surely, for every $\lambda\in\mathbb C$, there exists a geometric sequence such that
\[
x_{i+1}=\lambda x_i.
\]

Hence every $\lambda\in\mathbb C$ is an eigenvalue of $L$.

The eigenspace $V^L_\lambda$ associated with $\lambda$ is all sequences $(x_n)_{n=0}^{\infty}$ such that
\[
x_{i+1}=\lambda x_i
\]
for all $i=0,\ldots,\infty$.

If $x_0\neq0$, then
\[
x_n=\lambda^n x_0.
\]

So
\[
V^L_\lambda=\Span\{(1,\lambda,\lambda^2,\ldots)\}.
\]

\subsection*{(ii)}

\[
R(x_0,x_1,\ldots)=(0,x_0,x_1,\ldots).
\]

If $\lambda$ is an eigenvalue of $R$, then
\[
R(x_0,x_1,\ldots)=\lambda(x_0,x_1,\ldots)
\]
\[
=(0,x_0,x_1,\ldots).
\]

Entrywise,
\[
\lambda x_0=0,
\]
\[
\lambda x_1=x_0,
\]
\[
\lambda x_2=x_1.
\]

So
\[
\lambda x_i=x_{i-1}
\]
for all $i=1,\ldots,\infty$.

Let $x_{-1}=0$.

Then let's consider two cases.

Case 1: $\lambda=0$. Then
\[
x=(0,\ldots,0),
\]
which contradicts that $\lambda$ is an eigenvalue of $R$.

Case 2: $\lambda\neq0$.

Then
\[
\lambda x_0=0\Rightarrow x_0=0.
\]

Consequently,
\[
\lambda x_1=0\Rightarrow x_1=0,
\]
and generally
\[
\lambda x_i=0
\]
for all $i=0,\ldots,\infty$.

Hence
\[
x_i=0
\]
for all $i=0,\ldots,\infty$, which contradicts that $x$ is an eigenvector.

Therefore, by two cases,
\[
\spec(R)=\varnothing.
\]

\end{document}
